% !TeX root = ../papers/nc-p3.tex
\section{Walsh族を通る二つの大域的な半周期成分}
\label{sec:half-period-components}
\begin{theorem}\label{thm:half-period-components}
十五次元成分の安定商で、一般Walsh族を含む二位置移動の被約不動軌跡の既約成分は
九次元と五次元の二つである。前者の一般点の射影安定化群は自明、後者では位数2である。
後者の三つの局所枝はWalsh族に沿って全対称群 $S_3$ で置換される。
その正規化は一般Walsh点の上に三点を持ち、各枝はそこで滑らかである。
前者の正規化は一般Walsh点の上に一点を持つが、その近くに正規化の二点が同定される
五次元の部分が三つある。
\end{theorem}
\begin{proof}
前節の線型切片を使う。$x=0$ の九次元空間の一つを $L$ とし、六対の交換符号を全て正に取る。
この空間を保つ群は $D=\{(0,z)\}$ である。
各 $h\ne0$ の二つの独立座標に、同じ $(-1)^{z\cdot h}$ が働く。
全座標が非零なら、二位置移動を含まない安定化群は自明である。
$L$ の軌道の有限商の正規化は $L/D$ となる。これは有限な双有理写像で、有限線型群の商は正規である。
$G$ のうち $L$ を保たない元が $L$ の点を再び $L$ に送るには、その点の
非自明な文字の支持がただ一つの $h$ に含まれる必要がある。
実際、移動部分 $y\ne0$ は $y\cdot h^\flat\ne0$ の対の符号比を反転する。
$h=y$ の二座標を $u,v$ とすると、$\{(y,z)\}$ は同時符号のほかに相対符号を反転する。
よって $uv\ne0$ では $L/D$ の二点が像で同一となる。各 $h$ につきこの部分は
三自明方向と二座標の五次元である。六座標が零のWalsh点上には正規化の一点だけがある。

$x\ne0$ の局所枝では、$h=x$ の二対だけを使い一方の符号を逆にする。
両対が非零なら実際の安定化群は $\{1,(0,x^\flat)\}$ である。
持ち上げた二位置移動 $g\tau$ の二乗はその非自明元で、持上げの射影位数は4である。
固定した $x$ の一方の線型空間を保つ群は移動成分が $0,x$ の元である。
二座標への像は二符号の独立な変更全体なので、その商は三自明座標と $u^2,v^2$ の滑らかな空間になる。
他方の共役空間との交わりは三自明座標のみである。二つの線型空間の関数環は
その交叉上で一致する関数の対であり、交叉に働く共役作用は恒等であるから、
両者を合わせて $G$ で割ってもこの一つの滑らかな商を得る。
異なる三つの $x$ の支持は互いに異なり、交叉はWalsh方向である。

大域的には、欠色を固定した既約Walsh表示上で一定の $g\tau$ を固定する。
一般の表示点の近くでは滑らかな周囲空間への有限位数自己同型の不動点は滑らかである。
従って一般表示を含むこの固定点の既約成分は一つで、その基底変更飽和の閉包は既約である。
四基底を動かした場合も新しい局所成分は増えない。切片での有限個の固定空間に、
軌道内の $g\tau$ 不変方向を加えれば表示の固定スキームの接空間であり、
飽和して商に移せば対応する線型切片の全次元が出る。
Walsh表示を保つ基底変換は定理\ref{thm:walsh-presentation-space}で
移動ラベルの全 $\mathrm{GL}_2(\mathbb F_2)$ を実現する。
そのような変換 $n$ は $\tau$ を $g'\tau$ に共役し、
$g'$ は一般Walsh表示を固定するので $G$ に属する。
共役後も位数2であることから $g'$ の移動部分は零である。
ゆえに三つの非零 $x$ は全置換で移る。各 $x$ の二つの符号も既に $G$ で共役なので、
五次元の全ての局所枝の大域的な飽和は一致する。九次元側も同様に一成分であり、
二つの次元と局所非包含から両者は不動軌跡の別の既約成分である。
三つの標識は一般Walsh表示上で一定で、その指定を忘れる有限被覆の変換が
全六置換を与えるので、五次元側の正規化の三点の一般モノドロミーは $S_3$ である。
\end{proof}

\begin{proposition}\label{prop:half-period-pencils}
五次元成分の一般点では、二位置移動の位数4の持上げを取ると、各隣接切断は
八次の対称行列の階数4部分と八次の交代行列の階数2部分に分かれる。
交代部分の二つの対角ブロックの像は $2\times2$ 行列の二平面で、各々に
階数1の線が相異なって二本ある。これらは全頂点の四直線の基底枠を指定なしで有限通りに決める。
九次元成分の一般点では二位置移動の持上げは位数2で、両隣接切断を対称行列にできる。
\end{proposition}
\begin{proof}
二位置移動で対向する空間を同定する。五次元枝で持上げの二乗は二対二の固有空間を持つ
二つの射影対合である。偶、奇の符号を $p_i,q_j$ とすると、第一切断 $M$ は
$M_{ij,kl}=p_i q_j M_{kl,ij}$、$p_i q_j=p_k q_l$ に台を持つ形にできる。
第二切断は $N_{kj,il}=p_i M_{ij,kl}$ とすれば、同じ性質を $p_k q_j$ について持つ。
この符号の型は基点の実際の持上げを以下で確認すれば近傍でも一定である。

Walsh点で持上げを $g\tau$、$z=0$ とし、偶・奇四次行列を $A,B$ とする。
文字 $i$ の列をそれぞれ $i+x$ へ運び、$J(i)(-1)^{x_1 i_0},J(i)$ を掛ける。
第一切断を $M=W(A\otimes B)^{\mathsf t}$ に取れば
$A^2$ の文字は $x^\flat$、$B^2$ の文字も共通倍率 $J(x)$ を除いて $x^\flat$ である。
よって反対称側に来る非欠色 $b$ は各非零 $x$ で一つだけであり、切断の階数は順に4,2である。
階数の和は6なので、この二下界は近傍で等号のまま保たれる。
二つの負符号ブロックの内部の四次交代行列について、
\path{computations/verify_half_period_components.py} は先の53の点で
各像が二次元、射影鉛筆の行列式に二単根があることを三つの持上げ全てで確かめる。
単根の二行列の像と核はそれぞれ独立な二線である。
負符号空間は $E_+\otimes F_-$ と $E_-\otimes F_+$ の直和なので、これで全四空間の二分解を
一度ずつさらに二分できる。一般点では安定化群が位数2で、二つの持上げは互いに逆であるため
これらを両方取っても有限個の枠しか生じない。九次元の $x=0$ の場合は
$A,B$ が対角の符号行列で、双方の二乗とも1である。
同じ同定で全体が対称となり、族の対応する固定スキームでこの形が保たれる。
\end{proof}
